The target may be either be specified by \textbf{regions} of homogeneous 
material, or by quasi-continually varying atom densities discretized in cells 
of a regular grid.

A region is defined by its material and geometry. The \textbf{material} is
primarily defined by its name, which is a chemical formula similar to the ion
name. The \textbf{density} of a region is set to the weighted densities of the
elementary materials by default, or can be specified by the user. Note that for
compound targets, the weighted densities are often inaccurate, so their
densities should be specified in the input file. 

Another material property is its \textbf{crystallinity}. If the material is
crystalline, the crystal structure is defined in the \texttt{CRYSTAL} record. In
addition, the crystallographic directions of the surface normal and of a surface
vector may be specified. Alternatively, the material may be specified as
polycrystalline, in which case a new random orientation of the crystal is chosen
for each ion. 

A further material property is the RMS \textbf{thermal vibration} amplitude of
the atoms, which may either be calculated from the Debye theory using the target
and Debye temperature, or it may be specified directly. Thermal vibrations are
only relevant in crystalline targets.

When the target is defined by the atom densities in cells of a regular
grid, the material properties in each cell are then linearly interpolated
from the properties of the elementary materials. The cell-based definition
of the material allows for spatially dependent material properties (within
the resolution of the grid). It may be used both for static and dynamic 
(see Section \ref{s:updates}) simulations. The use of cell-based material
specification is currently restricted to 1-D geometries.

The region-based material definition can also be used both for static and
dynamic simulations. For dynamic simulations, the materials and regions 
specifed on the \texttt{\&MATERIAL} and \texttt{\&GEOMETRY} record are
converted into a cell-based description during the intialization phase.
The cell densities are then allowed to change during the simulation. Since the
other material properties depend on the target composition, they can only be 
interpolated from the elemental material properties. Crystalline targets are not
allowed in dynamic simulations. 


\begin{center}
\begin{tabular}{lll}
parameter \quad                 & IMSIL name    & to be specified in record \\
\hline
number of regions               & \texttt{NR}        & \texttt{\&SETUP} \\
cell-based defintion flag       & \texttt{USECELL}   & \texttt{\&SETUP} \\
file containing cell densities  & \texttt{CELLFILE}  & \texttt{\&SETUP} \\
material name                   & \texttt{NAME}      & \texttt{\&MATERIAL} \\
elementary material density     & \texttt{DENSITY}   & \texttt{\&ATOMS} \\
material density                & \texttt{DENSITY}   & \texttt{\&MATERIAL} \\
crystallinity flag              & \texttt{XTAL}      & \texttt{\&MATERIAL} \\
polycrystallinity flag          & \texttt{LPOLY}     & \texttt{\&MATERIAL} \\
surface normal                  & \texttt{WAFER}     & \texttt{\&MATERIAL} \\
surface vector                  & \texttt{VSURF}     & \texttt{\&MATERIAL} \\
target temperature              & \texttt{TEMP}      & \texttt{\&SETUP} \\
Debye temperature               & \texttt{TDEBYE}    & \texttt{\&MATERIAL} \\
RMS thermal vibration amplitude & \texttt{VIB}       & \texttt{\&MATERIAL} \\

\end{tabular}
\end{center}

